\section{Synthetic data: RFNN spectral microscope}
\label{sec:rfnn}

\subsection{Setup and assumptions}
\label{sec:rfnn-setup}

We adopt the RFNN framework of \citet{bonnaire2025} unchanged: a two-layer network
\begin{equation}
  s_A(x_t) \;=\; \tfrac{1}{\sqrt{\pwidth}}\,A\,
  \tanh\!\bigl(\tfrac{1}{\sqrt{\dlat}}\,W x_t\bigr),
  \label{eq:rfnn}
\end{equation}
with frozen first-layer weights $W \in \R^{\pwidth \times \dlat}$ drawn i.i.d.\ from $\mathcal{N}(0, 1)$, $\tanh$ activation, and a trained readout $A \in \R^{\dlat \times \pwidth}$ initialized to zero. Training is score matching at a single fixed diffusion time $t = 0.01$, and the analysis is a proportional-limit statement ($\dlat,\nsamp,\pwidth\to\infty$ at fixed ratios), but our sweeps hold $\nsamp=500$ fixed while $\dlat$ runs from $5$ to $200$, so $\psi_n=\nsamp/\dlat$ moves from $100$ to $2.5$; all finite-size statements below are therefore at $O(\dlat^{-1/2})$ accuracy and nothing is uniform along the sweep (Sec.~\ref{app:theory-scope}). We use three width choices to separate the spectral mechanism from width artifacts; Appendix~\ref{app:width-controls} spells out these controls, and Appendix~\ref{app:rfnn-clean} reports the complete sweeps. First, $\pwidth=64\dlat$ preserves the original RFNN scaling from \citet{bonnaire2025}: the random feature map stays highly overcomplete as $\dlat$ grows, so this setting asks whether the four-bulk ordering is visible in the well-resolved, many-feature regime (Figures~\ref{fig:app-rfnn-exp2-widths} and \ref{fig:app-rfnn-exp3-widths}). Second, $\pwidth=\dlat+\nsamp+300$ grows only enough to represent the $\nsamp$ non-null feature directions plus a margin of $\dlat+300$ (the rank-null block has $\pwidth-\nsamp$ modes, not $\pwidth-\dlat-\nsamp$); here $\psi_p=\pwidth/\dlat$ falls toward $1$ and the noise-dimension block broadens by $((1+\sqrt{D/\pwidth})/(1-\sqrt{D/\pwidth}))^2$, so four-block separation is predicted to fail for $\dlat\gtrsim100$ at $\signoise=0.5$ (Rem.~\ref{rem:scope-psip}). This controls for the possibility that the separation is an artifact of the much wider $64\dlat$ feature map (Figures~\ref{fig:app-rfnn-exp2-pndr} and \ref{fig:app-rfnn-exp3-pndr}). Third, the fixed-$\pwidth$ ablation uses fixed total null energy, rescaling $\sigma_\perp^2\propto 1/(\dlat-\dint)$ in the latent-dimension sweep. This is the strictest control: increasing $\dlat$ no longer increases either RFNN width or total injected null variance, so any remaining separation is tied to how the same null energy is spread across more latent coordinates; the theory predicts that for $\dlat-\dint\gg 50c$ this control's spectrum coincides with the $\signoise\to0$ spectrum (noise-dimension block on the diffusion floor, three blocks), with $\taugen$ still growing linearly in $\dlat$ (Prop.~\ref{prop:scope-control}) (Figures~\ref{fig:app-rfnn-exp2-pfixed} and \ref{fig:app-rfnn-exp3-pfixed}). Because all width controls support the same qualitative conclusion, the main text shows the strictest representative setting and the remaining RFNN sweeps are reported in Appendix~\ref{app:rfnn-clean}.

\textbf{Training dynamics.}
Freezing $W$ makes the score-matching loss quadratic in $A$, so full-batch gradient flow on $A$ is the linear ODE
\begin{equation}
  \tfrac{\mathrm{d}}{\mathrm{d} T} A(T) \;=\; -\bigl(A(T) - A^{\star}\bigr)\,\Umat,
  \label{eq:rfnn-gradflow}
\end{equation}
where $A^\star$ is the train-loss minimizer and the empirical feature correlation matrix is
\begin{align}
  \Umat \;&=\; \tfrac{1}{\nsamp}\sum_{\mu=1}^{\nsamp}
    \E_{\xi}\!\bigl[\phi(x_t^{\mu})\phi(x_t^{\mu})^{\top}\bigr],
  \label{eq:U} \\
  \phi(x_t) \;&:=\; \tfrac{1}{\sqrt{\pwidth}}\,
    \tanh\!\bigl(\tfrac{1}{\sqrt{\dlat}}\, W x_t\bigr),
\end{align}
with $x_t^\mu$ the noisy diffused state of training point $\mu$ and $\xi$ the diffusion noise. Diagonalizing Eq.~\eqref{eq:rfnn-gradflow} in the eigenbasis of $\Umat$ gives an exponential per-mode evolution
\begin{equation}
  a_i(T) - a_i^{\star} \;=\; \bigl(a_i(0) - a_i^{\star}\bigr)\,
    e^{-\lambda_i T},
  \label{eq:rfnn-mode-decay}
\end{equation}
so eigenmode $i$ of $\Umat$ is absorbed by the readout with characteristic timescale $\tau_i=1/\lambda_i$ in the flow time $T$ of Eq.~\eqref{eq:rfnn-gradflow}; with the code's learning rate $0.01\dlat/\Delta_t$ and its loss normalization, $T=0.02\,\dlat\times$steps, and the per-mode \emph{risk} decays twice as fast as the parameter deviation (Rem.~\ref{rem:clock-p1}, Rem.~\ref{rem:lin-gd}). All modes decay in parallel: no mode waits for another.

\textbf{Estimating $\Umat$.}
$\Umat$ is computed numerically per run: we fix the training set and the random first layer $W$, draw a Monte Carlo estimate of the expectation over diffusion noise $\xi$ in Eq.~\eqref{eq:U} (50~Monte Carlo noise samples for the released $\pwidth=64\dlat$ spectra, \texttt{n\_noise\_samples=50} in the run configs; 500 for the $\pwidth=\dlat+\nsamp+300$ set; Appendix~\ref{app:methods}). Because the expectation over $\xi$ is only Monte-Carlo approximated, the reported $\Umat$ has full rank $\pwidth$ with a soft shoulder at index $\nsamp$ (eigenvalue ratio $1.0$--$1.7$); the exact-rank cliff at $\nsamp$ holds for the single-draw matrix (\cref{rem:fourbulk-singlecopy}), assemble the $\pwidth \times \pwidth$ matrix, and then take its eigendecomposition. The reported spectra are the eigenvalues of this empirical estimate; no analytic surrogate is used.

\textbf{What the eigenvalues mean.}
Eq.~\eqref{eq:rfnn-mode-decay} says each $\lambda_i$ is a learning rate for one feature direction: large $\lambda_i$ are absorbed quickly, small ones slowly, and sorting the spectrum sorts these rates; nothing waits for anything else. The RFNN is therefore performing weighted PCA on a frozen nonlinear basis: the random map $\phi$ defines a fixed set of candidate directions in $\R^{\pwidth}$, $\Umat$ is the empirical covariance of the training data after that lift, and $\lambda_i$ measures how strongly the training samples vote for direction $i$. Modes are learned in order of how well the data supports them. Our analysis operates entirely inside this framework; only the data covariance $\Sigmadata$ changes, and Bonnaire et al.'s theorem already allows an arbitrary spectral measure of $\Sigmadata$: what we add is the evaluation at the two-atom measure with $\signoise>0$ and the training-time reading of the result.

\subsection{From two bulks to four}
\label{sec:rfnn-fourbulks}

\citeauthor{bonnaire2025} prove (Theorem 3.1, arbitrary spectral measure of $\Sigmadata$) a Gaussian-equivalent characterization of the spectrum of $\Umat$, which for isotropic $\Sigmadata=\mathbb I$ splits into a population bulk $\rho_2$ and a sample bulk $\rho_1$. With anisotropic data covariance $\Sigmadata = \diag\!\bigl(\sigsig^2 \mathbb{I}_{\dint},\, \signoise^2 \mathbb{I}_{\dlat - \dint}\bigr)$ ($\sigsig \gg \signoise$), the sorted spectrum splits into \emph{four} index blocks (\cref{prop:fourbulk-main}): (i) $\dint$ signal modes, a finite-rank spike set at $a_1(q)^2\alpha_t^2/\dlat$ (BBP-shifted by the null block); (ii) $\dlat-\dint$ noise-dimension modes filling a Marchenko--Pastur band around $a_1(q)^2\beta_t^2/\dlat$; (iii) $\nsamp-\dlat$ sample-specific modes, the residual (nonlinear) part of the features with total mass $a_*(q)^2$ and mean $\approx a_*(q)^2/\nsamp$; (iv) $\pwidth-\nsamp$ rank-null modes. Blocks (i)--(ii) are Bonnaire's $\rho_2$ and block (iii) is their $\rho_1$; the sample block does \emph{not} split into signal and null sub-blocks. We observe these four populations in every configuration we ran (Figure~\ref{fig:fourbulk-clean}; complete Experiment~1 width controls in Appendix Figures~\ref{fig:app-rfnn-exp2-widths}--\ref{fig:app-rfnn-exp2-pfixed}; Experiment~2 controls in Figures~\ref{fig:app-rfnn-exp3-widths}--\ref{fig:app-rfnn-exp3-pfixed}; Gaussian robustness checks in Figures~\ref{fig:app-rfnn-robustness}-- \ref{fig:app-rfnn-robust-dint-sigscale}).

\begin{figure*}[!t]
  \centering
\includegraphics[width=0.86\textwidth,height=0.62\textheight,keepaspectratio]{figures/rfnn_exp2_pndr_main_selected.pdf}
  \caption{\textbf{Experiment 1: RFNN four-bulk spectrum for the
    $\dlat$ sweep.}
    Representative panels from the full latent-dimension sweep are shown here for readability; the complete sweep appears in Appendix Figure~\ref{fig:app-rfnn-exp2-pndr}. A synthetic GMM with fixed $\dint=5$ is embedded into increasing $\dlat$. The red signal modes stay fixed in count, while the green noise-dimension block grows in count with $\dlat-\dint$ and the blue sample-specific block shrinks in count as $\nsamp-\dlat$; the gap between the smallest red eigenvalue and the largest blue eigenvalue widens with $\dlat$. Because $\dint$ is known exactly in this synthetic construction, the regime change is also exact: at $\dlat=\dint$ there is no green buffer, and increasing $\dlat$ increases the buffer almost purely monotonically rather than producing irregular fitted subclusters. Colors are assigned by sorted eigenvalue-index blocks rather than fitted clusters: red is signal/source, green is latent-null/noise-dimension, blue is sample-specific, and purple is rank-null; dashed vertical lines mark within-block median log-eigenvalues.}
  \label{fig:fourbulk-clean}
\end{figure*}

The coloring in these plots is deliberately non-adaptive. For each run we sort the eigenvalues of $\Umat$ in decreasing order and then assign labels by the \emph{a priori} index blocks predicted by the data construction: indices $1{:}\dint$ are signal/source, indices $\dint{+}1{:}\dlat$ are latent-null or noise-dimensional, indices $\dlat{+}1{:}\nsamp$ are sample-specific, and indices $\nsamp{+}1{:}\pwidth$ are rank-null (for the single-draw $\Umat$ these are exact zeros; for the noise-averaged $\Umat$ we report they form a residual tail of mass $\lesssim 20\%$ of $a_*(q)^2$). At $\dlat=5$, $\pwidth=320<\nsamp$ and there is no rank-null block. In the over-intrinsic sweep, the signal/source block is clipped at the observed dimension $\min(\dint,\dlat)$. No clustering or threshold fitting is used to place the colors on the histogram. Thus, when the colored groups separate on the $x$-axis, the claim is stronger than a post-hoc visual partition: the eigenvalue order itself has organized into the predicted latent, sample, and rank blocks. The gap at index $\dint$ between signal and noise-dimension modes is $O(\alpha_t^2/\beta_t^2)$ at both noise scales; the gap at $\dlat$ between noise-dimension and sample modes is $0.6$--$1.3$ decades at $\signoise=0.5$ and closes ($\lambda_{\dlat}/\lambda_{\dlat+1}=1.0$--$1.15$) at $\signoise=0.01$, where the noise-dimension block sits on the diffusion floor $a_1(q)^2\Delta_t/\dlat$ and merges with the sample block (Prop.~\ref{prop:scope-sigperp}).

\citet{bonnaire2025} already note (Figure~4, right panel) that $\rho_2$ develops internal structure when $\Sigmadata$ has multiple eigenvalues, but at signal-to-null ratio $3{\times}$ they treat it as a single deformed bulk. The latent-diffusion regime sits at the opposite extreme ($\sigsig^2/\signoise^2 \approx 10^4$ at $\signoise = 0.01$, $\approx 4$ at $\signoise = 0.5$): the signal spikes and the noise-dimension band separate by the factor $\alpha_t^2/\beta_t^2$ ($\approx10$ at $\signoise=0.5$, $\approx140$ at $\signoise=0.01$ for $t=0.01$), which is why it is useful to count them as distinct and analyze their properties separately.

The complementary source-dimensionality sweep asks the reverse question from the $\dlat$ sweep. Instead of making the observed latent space wider around a fixed five-dimensional source, we hold the observed latent space fixed at $\dlat=20$ and regenerate the source distribution with increasing $\dint$. If the mechanism is really an excess-observed-dimension buffer, then the green latent-null block should shrink as $\dint$ approaches $\dlat$, and it should vanish once all observed coordinates are needed to carry source variation. The over-intrinsic cases $\dint>\dlat$ are a compression stress test: adding source dimensions beyond the observed latent dimension cannot create extra observed signal modes, so the RFNN should no longer contain a separate null buffer. This is Experiment~2 in the appendix, where we repeat the sweep under the same three width controls (Figures~\ref{fig:app-rfnn-exp3-widths}--\ref{fig:app-rfnn-exp3-pfixed}).

\textbf{Mechanism: a noise-dimension buffer.}
\looseness -1
Every mode decays in parallel with rate $\lambda_i$, so there is no sequential traversal: $\taugen=1/\lambda^{\rm sig}_{\min}$ and the onset of memorization $\taumem^{\rm on}=1/\lambda^{\rm samp}_{\max}$ are set by two edges of the spectrum, and the buffer is their ratio $R=\lambda^{\rm sig}_{\min}/\lambda^{\rm samp}_{\max}$ (Cor.~\ref{cor:buffer-buf}). On the released spectra $\taugen$ grows $11\times$ in flow units over $\dlat=5\to200$ (it is not $\dlat$-independent; in code steps it \emph{falls} because the learning rate scales with $\dlat$), $\taumem^{\rm on}$ grows $220\times$, and $R$ grows monotonically from $21$ to $430$. The growth comes from the fall of $q=\operatorname{tr}M_t/\dlat$ ($2.76\to0.33$), which raises $a_1(q)^2$ ($0.18\to0.63$) and lowers $a_*(q)^2$ ($0.079\to0.0047$), and from the collapse of the small-$\dlat$ excess of the sample-block top over its Marchenko--Pastur edge. The count $\dlat-\dint$ enters no timescale. In the isotropic case the $\dlat-\dint$ noise-dimension modes still exist but are not spectrally separated from the signal modes; the four-block \emph{structure} is absent, and $R$ is predicted to rise to a maximum near $\dlat\approx50$ and then decline like $\nsamp/\dlat$ (Rem.~\ref{rem:isotropic-buf}).

\textbf{What the gen-gap measures.}
\looseness -1
The training set carries two kinds of structure. The data bulks correspond to \emph{population} structure: directions of variance that exist in the population and happen to also be visible in the training sample. The sample bulk corresponds to \emph{null structure}: directions that look like signal in the $\nsamp$ training points but do not exist in the population. The RFNN absorbs eigenmodes in order of decreasing $\lambda_i$, so it fits the population structure first and the null structure later. While it is fitting population structure, train and test loss drop together. Once it begins fitting null structure, train loss keeps dropping but test loss does not, and the gap opens. The gen-gap is a measure of how much sample-specific structure the readout has absorbed (Prop.~\ref{prop:gap-p1}); it opens at the sample-block edge, and widening $\dlat$ delays that edge through the fall of $a_*(q)^2$. The logged proxy \texttt{gen\_gap}$>0.02$ is contaminated by a single-draw noise floor $\sqrt{2/(\nsamp\dlat)}$ that exceeds the threshold at small $\dlat$, so onset times should be re-extracted with a relative threshold (Rem.~\ref{rem:proxy-p1}).

\textbf{Decomposing the four-bulk structure.}
The signal and noise-dimension blocks of $\Umat$ inherit the block structure of $\Sigmadata$ directly: $X^\top X / \nsamp$ on the same training data has $\dint$ large eigenvalues in the signal subspace and $\dlat - \dint$ small eigenvalues in the null directions, separated by $\alpha_t^2/\beta_t^2$, and the lift by $W$ carries them into the data blocks of $\Umat$ multiplied by $a_1(q)^2/\dlat$, where $a_1(q)=\E[\tanh'(\sqrt q z)]$ depends on $\dlat$ through $q$ (Lemma~\ref{lem:q-p1}); this factor, not a constant, is what makes every edge $\dlat$-dependent. The sample block is not an artifact: it is the residual, nonlinear part of the features ($a_*(q)\,z$ in the Gaussian-equivalent decomposition), carries total mass $a_*(q)^2$, and is exactly the part of the model that can encode individual training points; the rank-null block ($\pwidth-\nsamp$ modes) is the genuine feature-space excess. The buffer mechanism is gradient flow following the data's variance hierarchy \emph{after} the nonlinearity has redistributed it between the linear and residual parts.


